\documentclass[12pt,a4paper]{article}

\usepackage[utf8]{inputenc}
\usepackage[spanish]{babel}
\usepackage{geometry}
\usepackage{hyperref}
\usepackage{xcolor}
\usepackage{booktabs}
\usepackage{tabularx}
\usepackage{parskip}

\geometry{top=2.5cm,bottom=2.5cm,left=2.5cm,right=2.5cm}
\hypersetup{colorlinks=true,urlcolor=blue}
\pagestyle{empty}

\begin{document}

\begin{center}
{\Large\bfseries Informe de Reflexi\'on --- Examen Final CI/CD}\\[0.2cm]
{\small nnyez --- 29 de julio de 2026 --- \url{https://github.com/nnyez/inventario-app}}
\end{center}

\subsection*{1. Justificación de la estrategia Canary}

Se eligió Canary sobre Blue-Green porque Kubernetes distribuye tráfico proporcional a la cantidad de pods que matchean el selector del Service. Con 4 réplicas estables (v1/blue) y 1 canary (v2/green), el $\sim80\%$ del tráfico va a la versión estable y $\sim20\%$ a la nueva. Canary requiere solo 5 pods frente a 8 de Blue-Green, es más eficiente para una aplicación liviana como \texttt{inventario-app}, y permite rollback inmediato escalando el canary a 0 réplicas. Además, permite aumentar gradualmente el porcentaje de tráfico a la nueva versión simplemente escalando el número de réplicas canary.

\subsection*{2. Observación sobre persistencia de datos}

Se cre\'o un producto vía API (4 productos total) y se forz\'o la eliminaci\'on de un pod. Tras la recreaci\'on, el cat\'alogo conten\'ia solo 3 productos (los seed originales). Esto ocurre porque la base de datos es un archivo JSON local (\texttt{data/products.json}) en el filesystem ef\'imero del pod. No es un error, es el comportamiento esperado sin PVC o base de datos externa.

\subsection*{3. Problemas reales encontrados}

\small
\begin{tabularx}{\textwidth}{p{3cm}p{3.5cm}X}
\toprule
\textbf{Problema} & \textbf{Causa} & \textbf{Solución} \\
\midrule
Pipeline fallaba & \texttt{trivy-action@0.28.0} no existe & Cambiar a @v0.36.0 \\
Build fallaba por caché & \texttt{cache-to=type=gha} incompatible & Eliminar caché \\
Error ``denied: not allowed'' & GITHUB\_TOKEN sin packages:write & Agregar permissions \\
Trivy CRITICALs & node:18-alpine con CVE OpenSSL/tar & Alpine 3.21, apk upgrade \\
Pods no arrancaban & Secret no exist\'ia & kubectl create secret \\
Sin balanceo port-forward & port-forward no distribuye tráfico & Probar dentro del clúster \\
\bottomrule
\end{tabularx}
\normalsize

\subsection*{4. Métricas DORA}

Se calcularon con datos verificables de tres fuentes: \texttt{git log} (timestamp de commits), \texttt{gh run list} (timestamp de pipelines), y \texttt{kubectl describe deployment} (campo \texttt{restartedAt} para el momento exacto en que el cambio lleg\'o al cl\'uster).

\small
\begin{tabularx}{\textwidth}{Xcc}
\toprule
\textbf{Métrica} & \textbf{Valor} & \textbf{Nivel DORA} \\
\midrule
\textit{Lead Time} (commit $\rightarrow$ clúster): & 16 seg -- 41 min & \textbf{Élite} \\
Despliegue inicial: \texttt{3ab1521} 04:23:30 $\rightarrow$ 05:04:29 & 41 min & ($<$1 h) \\
Namespace+secrets: \texttt{112578a} 05:06:48 $\rightarrow$ 05:07:04 & 16 seg & \\
STARTUP\_DELAY: \texttt{2d25c12} 05:07:51 $\rightarrow$ $\sim$05:08 & $\sim$1 min & \\
\midrule
\textit{Deployment Frequency}: cambios promovidos al clúster & 5 deploys/día & \textbf{Alto} \\
\midrule
\textit{Change Failure Rate}: pipelines fallidos / total & 7/14 = 50\% & \textbf{Bajo} \\
\bottomrule
\end{tabularx}
\normalsize

\subsection*{5. Reflexión}

El CFR del 50\% refleja la etapa inicial de prueba y error (7 fallos por autenticación, versión incorrecta de acciones, caché de Docker y vulnerabilidades CRITICAL detectadas por Trivy). En un entorno maduro se espera reducirlo por debajo del 15\% (nivel Alto/Élite). El lead time \textbf{Élite} (16 segundos para cambios pequeños como configuración de namespace o STARTUP\_DELAY) demuestra que el pipeline entrega cambios casi instantáneamente cuando está estable. La frecuencia \textbf{Alta} (5 deploys/día) es propia de un desarrollo iterativo donde se experimenta y corrige sobre infraestructura real. En conjunto, las métricas confirman que las pr\'acticas implementadas (fail-fast entre jobs, escaneo de seguridad con Trivy, readiness probe tolerante al arranque lento, y manejo de secretos vía \texttt{secretKeyRef}) mejoran significativamente la calidad y confiabilidad del despliegue.

\end{document}
